% Generated draft: requires mathematical and source review.
\documentclass[11pt]{article}
\usepackage[margin=25mm]{geometry}
\usepackage{fontspec}
\setmainfont{Latin Modern Roman}
\usepackage{amsmath,amssymb,mathtools,unicode-math}
\setmathfont{Latin Modern Math}
\usepackage{enumitem,needspace,xurl,hyperref}
\setlength{\parindent}{0pt}
\setlength{\parskip}{5pt}
\setlength{\emergencystretch}{3em}
\newcommand{\N}{\mathbb{N}}
\newcommand{\Z}{\mathbb{Z}}
\newcommand{\Q}{\mathbb{Q}}
\newcommand{\R}{\mathbb{R}}
\newcommand{\C}{\mathbb{C}}
\newcommand{\F}{\mathbb{F}}
\newcommand{\A}{\mathbb{A}}
\newcommand{\PP}{\mathbb{P}}
\newcommand{\E}{\mathbb{E}}
\newcommand{\eps}{\varepsilon}
\newcommand{\dd}{\,\mathrm{d}}
\begin{document}
\section*{UnsolvedMath ID: 30000310}
\textbf{Source problem code:} catalogue entry
\subsection{Definitions and mathematical statement}
Let \(u=(u_{ij})\) be the corepresentation in the source's universal Hopf algebra. Its character is
\[
\chi(u)=\sum_i u_{ii}.
\]
In a usual Hopf \(*\)-algebra or compact-quantum-group setting the corepresentation relation is \(\Delta(u_{ij})=\sum_k u_{ik}\otimes u_{kj}\), but the particular universal relations, dimension, completion, and hypotheses are absent from the record. The source asks for finer versions of its Theorem 2 with \(\chi(u)\) replaced by arbitrary coefficients of \(u\), such as individual \(u_{ij}\), or source-allowed matrix coefficients relative to specified vectors. The formulas and even the precise scope of Theorem 2 are not reproduced.

A spectral-measure question requires both an operator-algebra realization and a state. If \(x=x^*\) is a coefficient in a \(C^*\)-algebra \(A\) with a chosen state \(h\), its state spectral measure is the probability measure \(\mu_x\) determined by
\[
h(F(x))=\int_{\sigma(x)}F(t)\,\dd\mu_x(t)
\qquad(F\in C(\sigma(x))).
\]
For normal coefficients the analogous measure is on their complex spectrum. An arbitrary corepresentation coefficient need not be self-adjoint or normal; in that case a scalar spectral measure is not automatically defined by this formula. One can instead consider the noncommutative \(*\)-distribution given by mixed moments \(h(x^{\varepsilon_1}\cdots x^{\varepsilon_r})\), \(\varepsilon_j\in\{1,*\}\), but this is not silently substituted for the source's requested measure. A Haar state is a natural possible choice in a compact quantum group, yet the record does not identify its chosen state or assert that every intended universal Hopf algebra has such a realization.

The precise task is therefore to determine the distributions of the intended coefficients in the same sense and under the same hypotheses as the character results, after recovering those data. The supplied assessment checked 2026-08-21 calls this a natural finer question and reports it unresolved in the stated generality. It gives no explicit coefficient law to assert.
\subsection{Short English statement}
Extend the source's character-distribution formulas to individual corepresentation coefficients. The universal algebra, state, and meaning of spectral measure for possibly nonnormal coefficients must be specified before a precise distribution theorem is possible.
\subsection{Sources}
\begin{itemize}
\item Local UnsolvedMath registry, permanent ID 30000310, catalogue entry, \emph{Spectral Measures of Arbitrary Corepresentation Coefficients}.
\item \emph{C*-Algebren} (2006), Oberwolfach Reports, pp. 2305--2374, DOI \url{https://doi.org/10.4171/owr/2005/41}, Theorem 2 as referenced in the supplied original wording.
\item Supplied literature triage checked 2026-08-21 cites only that report. No universal relations, character formulas, coefficient law, or separate authoritative theorem citation is reproduced.
\end{itemize}
\end{document}
